\subsection{Identifying generalizable behavioral relationships}

As our motivating example illustrates, a traditional train-test split does not adequately guard against overfitting to a single data-generating process. 
This approach is statistically valid only when training and testing samples are independently drawn from the same distribution \citep{statisticalvapnik1998}. 
However, our objective differs fundamentally: we seek \persona{s} that remain predictive even when the underlying data-generating process shifts.
By ``data-generating process,'' we refer specifically to the experimental setting (the environment in which behavior occurs), the population from which behavior is observed, and the outcomes they generate.
In econometric terms, two processes diverge when the given input covariates appear with different distributions of values or when the covariates themselves are entirely different random variables.
For example, a model trained on \$20 dictator games may fail on \$5 dictator games due to different stake values (covariate shift) or when moving from dictator to public goods games (structural shift).

Without direct training data from the novel target setting, no standard statistical procedure ensures predictive accuracy \citep{ben2010theory,Klivans2024Domain}. 
Theoretically, only a fully specified causal model could guarantee accurate predictions \citep{Pearl2009Causality}. 
In practice, however, constructing or inferring such causal models generally requires strong, often unverifiable assumptions, which are rarely feasible in complex social science contexts.

Instead, we propose leveraging principles from invariant risk minimization \citep{arjovsky2020invariant} to identify behavioral relationships that remain stable despite shifts in the data-generating process. 
Rather than splitting a single dataset, we deliberately choose training data from one data-generating process (or several related processes) and validate using data from a related but distinct process from those used for training.
\Persona{s} that consistently predict behavior across these distinct datasets likely capture generalizable, and possibly causal \citep{Peters2016invariant, Heinze2018invariant}, drivers of behavior.
As a result, these validated \persona{s} should be more effective in predicting human responses in novel but theoretically similar settings.

Returning to our motivating example, suppose we have additional human data from a \$5 dictator game, with observed offers \{1,1,1,2,2,2\}.
Although stakes differ, both the \$20 and \$5 dictator games likely share a common decision-making process: individuals give slightly less than half the available amount, a reasonable balance more similar to what is observed in humans \citep{Henrich2001InSearch}.
By using the  \$20 dictator game as training data and the \$5 dictator game for testing, we are implicitly filtering for \persona{s} based on their capacity to predict this proportional response.
The earlier atheoretical prompt \emph{``You randomly choose numbers between 6 and 9,''} still fits the \$20 game perfectly, but clearly fails validation on the \$5 game.
By contrast, the theory-grounded \persona{} \emph{``You are self-interested but fair,''} likely generalizes effectively across both settings, and most importantly, the novel public goods game.

This validation process becomes even more robust when multiple distinct but theoretically-related datasets are available. 
Imagine dictator-game responses from \$1, \$5, \$10, \$20, \$50, and \$100 games, all exhibiting offers slightly below half.
Optimizing over multiple settings simultaneously further reduces the risk of overfitting. 
With each new training data-generating process, it is less likely that an arbitrary prompt will generalize in-sample.\footnote{We do not offer empirical examples of optimizing over multiple training settings in the main text (only single settings in Section~\ref{sec:predict_new_AR}).
Appendix~\ref{app:predict_new_CR} provides such analyses, including an additional preregistered experiment.}
Thus, if the underlying relationship governing these settings also holds for a novel target setting, the validated \persona{s} should robustly predict behavior there as well.

Yet, even with an effective validation method, a fundamental practical challenge remains: identifying the initial set of candidate \persona{s}. 
While our motivating example made this step appear straightforward, selecting plausible \persona{s} in more complex settings can be far less obvious.
In a perfect world, an optimization procedure would take in a massive set of prompts and filter them down to a smaller set that generalizes.
However, this does not yet work in practice because many different prompts can be used to generate the same behavior in a given setting.
For example, the number of possible natural language prompts that get an LLM to respond with numbers between 6 and 9 is enormous.
If we were to start from a large pool of prompts, then the optimization procedure might endlessly produce sets of \persona{s} that fit in-sample but fail on the validation set.
For now, it is essential to ground the candidate \persona{s} with a principled starting point, which---we argue and later demonstrate empirically---economic and behavioral theories offer.
